% Options for packages loaded elsewhere
\PassOptionsToPackage{unicode}{hyperref}
\PassOptionsToPackage{hyphens}{url}
\documentclass[
]{article}
\usepackage{xcolor}
\usepackage[margin=1in]{geometry}
\usepackage{amsmath,amssymb}
\setcounter{secnumdepth}{-\maxdimen} % remove section numbering
\usepackage{iftex}
\ifPDFTeX
  \usepackage[T1]{fontenc}
  \usepackage[utf8]{inputenc}
  \usepackage{textcomp} % provide euro and other symbols
\else % if luatex or xetex
  \usepackage{unicode-math} % this also loads fontspec
  \defaultfontfeatures{Scale=MatchLowercase}
  \defaultfontfeatures[\rmfamily]{Ligatures=TeX,Scale=1}
\fi
\usepackage{lmodern}
\ifPDFTeX\else
  % xetex/luatex font selection
\fi
% Use upquote if available, for straight quotes in verbatim environments
\IfFileExists{upquote.sty}{\usepackage{upquote}}{}
\IfFileExists{microtype.sty}{% use microtype if available
  \usepackage[]{microtype}
  \UseMicrotypeSet[protrusion]{basicmath} % disable protrusion for tt fonts
}{}
\makeatletter
\@ifundefined{KOMAClassName}{% if non-KOMA class
  \IfFileExists{parskip.sty}{%
    \usepackage{parskip}
  }{% else
    \setlength{\parindent}{0pt}
    \setlength{\parskip}{6pt plus 2pt minus 1pt}}
}{% if KOMA class
  \KOMAoptions{parskip=half}}
\makeatother
\usepackage{color}
\usepackage{fancyvrb}
\newcommand{\VerbBar}{|}
\newcommand{\VERB}{\Verb[commandchars=\\\{\}]}
\DefineVerbatimEnvironment{Highlighting}{Verbatim}{commandchars=\\\{\}}
% Add ',fontsize=\small' for more characters per line
\usepackage{framed}
\definecolor{shadecolor}{RGB}{248,248,248}
\newenvironment{Shaded}{\begin{snugshade}}{\end{snugshade}}
\newcommand{\AlertTok}[1]{\textcolor[rgb]{0.94,0.16,0.16}{#1}}
\newcommand{\AnnotationTok}[1]{\textcolor[rgb]{0.56,0.35,0.01}{\textbf{\textit{#1}}}}
\newcommand{\AttributeTok}[1]{\textcolor[rgb]{0.13,0.29,0.53}{#1}}
\newcommand{\BaseNTok}[1]{\textcolor[rgb]{0.00,0.00,0.81}{#1}}
\newcommand{\BuiltInTok}[1]{#1}
\newcommand{\CharTok}[1]{\textcolor[rgb]{0.31,0.60,0.02}{#1}}
\newcommand{\CommentTok}[1]{\textcolor[rgb]{0.56,0.35,0.01}{\textit{#1}}}
\newcommand{\CommentVarTok}[1]{\textcolor[rgb]{0.56,0.35,0.01}{\textbf{\textit{#1}}}}
\newcommand{\ConstantTok}[1]{\textcolor[rgb]{0.56,0.35,0.01}{#1}}
\newcommand{\ControlFlowTok}[1]{\textcolor[rgb]{0.13,0.29,0.53}{\textbf{#1}}}
\newcommand{\DataTypeTok}[1]{\textcolor[rgb]{0.13,0.29,0.53}{#1}}
\newcommand{\DecValTok}[1]{\textcolor[rgb]{0.00,0.00,0.81}{#1}}
\newcommand{\DocumentationTok}[1]{\textcolor[rgb]{0.56,0.35,0.01}{\textbf{\textit{#1}}}}
\newcommand{\ErrorTok}[1]{\textcolor[rgb]{0.64,0.00,0.00}{\textbf{#1}}}
\newcommand{\ExtensionTok}[1]{#1}
\newcommand{\FloatTok}[1]{\textcolor[rgb]{0.00,0.00,0.81}{#1}}
\newcommand{\FunctionTok}[1]{\textcolor[rgb]{0.13,0.29,0.53}{\textbf{#1}}}
\newcommand{\ImportTok}[1]{#1}
\newcommand{\InformationTok}[1]{\textcolor[rgb]{0.56,0.35,0.01}{\textbf{\textit{#1}}}}
\newcommand{\KeywordTok}[1]{\textcolor[rgb]{0.13,0.29,0.53}{\textbf{#1}}}
\newcommand{\NormalTok}[1]{#1}
\newcommand{\OperatorTok}[1]{\textcolor[rgb]{0.81,0.36,0.00}{\textbf{#1}}}
\newcommand{\OtherTok}[1]{\textcolor[rgb]{0.56,0.35,0.01}{#1}}
\newcommand{\PreprocessorTok}[1]{\textcolor[rgb]{0.56,0.35,0.01}{\textit{#1}}}
\newcommand{\RegionMarkerTok}[1]{#1}
\newcommand{\SpecialCharTok}[1]{\textcolor[rgb]{0.81,0.36,0.00}{\textbf{#1}}}
\newcommand{\SpecialStringTok}[1]{\textcolor[rgb]{0.31,0.60,0.02}{#1}}
\newcommand{\StringTok}[1]{\textcolor[rgb]{0.31,0.60,0.02}{#1}}
\newcommand{\VariableTok}[1]{\textcolor[rgb]{0.00,0.00,0.00}{#1}}
\newcommand{\VerbatimStringTok}[1]{\textcolor[rgb]{0.31,0.60,0.02}{#1}}
\newcommand{\WarningTok}[1]{\textcolor[rgb]{0.56,0.35,0.01}{\textbf{\textit{#1}}}}
\usepackage{longtable,booktabs,array}
\usepackage{calc} % for calculating minipage widths
% Correct order of tables after \paragraph or \subparagraph
\usepackage{etoolbox}
\makeatletter
\patchcmd\longtable{\par}{\if@noskipsec\mbox{}\fi\par}{}{}
\makeatother
% Allow footnotes in longtable head/foot
\IfFileExists{footnotehyper.sty}{\usepackage{footnotehyper}}{\usepackage{footnote}}
\makesavenoteenv{longtable}
\usepackage{graphicx}
\makeatletter
\newsavebox\pandoc@box
\newcommand*\pandocbounded[1]{% scales image to fit in text height/width
  \sbox\pandoc@box{#1}%
  \Gscale@div\@tempa{\textheight}{\dimexpr\ht\pandoc@box+\dp\pandoc@box\relax}%
  \Gscale@div\@tempb{\linewidth}{\wd\pandoc@box}%
  \ifdim\@tempb\p@<\@tempa\p@\let\@tempa\@tempb\fi% select the smaller of both
  \ifdim\@tempa\p@<\p@\scalebox{\@tempa}{\usebox\pandoc@box}%
  \else\usebox{\pandoc@box}%
  \fi%
}
% Set default figure placement to htbp
\def\fps@figure{htbp}
\makeatother
\setlength{\emergencystretch}{3em} % prevent overfull lines
\providecommand{\tightlist}{%
  \setlength{\itemsep}{0pt}\setlength{\parskip}{0pt}}
\usepackage{bookmark}
\IfFileExists{xurl.sty}{\usepackage{xurl}}{} % add URL line breaks if available
\urlstyle{same}
\hypersetup{
  pdftitle={Network modelling - Assignmen 1 - Group 24},
  pdfauthor={Andrin \ldots{}; Joshua \ldots{}; Lia \ldots{}; Stoyan Malinin},
  hidelinks,
  pdfcreator={LaTeX via pandoc}}

\title{Network modelling - Assignmen 1 - Group 24}
\author{Andrin \ldots{} \and Joshua \ldots{} \and Lia
\ldots{} \and Stoyan Malinin}
\date{}

\begin{document}
\maketitle

\section{Task 1}\label{task-1}

TODO

\section{Task 2}\label{task-2}

\textbf{TODO:} Think about whether to use one-sided or two sided tests
to report significance

\subsection{(1)}\label{section}

After running the QAP logistic regression, we see a slope coefficient of
2.93 with a p-value of 0.00. Therefore the result is significant and we
can state that friendship in W1 is positively correlated with friendship
in W2.

\subsection{(2)}\label{section-1}

The values are defined as follows:

\begin{Shaded}
\begin{Highlighting}[]
\NormalTok{literacyReceiver }\OtherTok{\textless{}{-}} \FunctionTok{matrix}\NormalTok{(attr}\SpecialCharTok{$}\NormalTok{literacy\_end, }\FunctionTok{length}\NormalTok{(ids), }\FunctionTok{length}\NormalTok{(ids), }\AttributeTok{byrow =} \ConstantTok{TRUE}\NormalTok{) }\CommentTok{\# (i)}
\NormalTok{sameGender }\OtherTok{\textless{}{-}} \FunctionTok{outer}\NormalTok{(attr}\SpecialCharTok{$}\NormalTok{gender, attr}\SpecialCharTok{$}\NormalTok{gender, }\StringTok{"=="}\NormalTok{) }\SpecialCharTok{*} \DecValTok{1} \CommentTok{\# (ii)}
\NormalTok{hiseiSender }\OtherTok{\textless{}{-}} \FunctionTok{matrix}\NormalTok{(attr}\SpecialCharTok{$}\NormalTok{HISEI, }\FunctionTok{length}\NormalTok{(ids), }\FunctionTok{length}\NormalTok{(ids), }\AttributeTok{byrow =} \ConstantTok{FALSE}\NormalTok{) }\CommentTok{\# (iii)}
\end{Highlighting}
\end{Shaded}

\begin{itemize}
\tightlist
\item
  For (i) we create a matrix where each column is filled with the
  literacy score of the corresponding student.
\item
  For (ii) we create a matrix with one in case of a gender match, and
  zero otherwise.
\item
  For (iii) we create a matrix where each row is filled with the HISEI
  score of each student
\end{itemize}

\subsection{(3)}\label{section-2}

Here are the coefficients for the regressors together with the
corresponding p-values

\begin{longtable}[]{@{}lll@{}}
\toprule\noalign{}
Name & Coefficient & p-value \\
\midrule\noalign{}
\endhead
\bottomrule\noalign{}
\endlastfoot
Friendship W1 & 2.318175 & 0.0000 \\
Literacy receiver & 0.0163386306 & 0.1282 \\
Same gender & 1.8688700 & 0.0000 \\
HISEI sender & 0.0290319336 & 0.1752 \\
\end{longtable}

Here are the corresponding interpretations (for the 5\% significance
level):

\begin{itemize}
\tightlist
\item
  The coefficient from (i) is insignificant, therefore we cannot confirm
  that higher literacy scores are associated with more friendship
  nominations received.
\item
  The coefficient from (ii) is positive and significant, so we can
  confirm the statement that having the same gender increases the
  likelihood of having a friendship nomination (while everything else
  stays the same).
\item
  The coefficient from (iii) is insignificant, so we cannot confirm the
  statement that higher HISEI values are associated with more friendship
  nominations send.
\end{itemize}

\emph{Note: Insignificant coefficients mean that we cannot reject the
null hypothesis that there is no correlation (the coefficient is 0).}

\subsection{(4)}\label{section-3}

\pandocbounded{\includegraphics[keepaspectratio]{solution_files/figure-latex/unnamed-chunk-2-1.pdf}}
We can here see the values for all coefficient across the different
classrooms.

\subsection{(5)}\label{section-4}

\begin{itemize}
\tightlist
\item
  \begin{enumerate}
  \def\labelenumi{(\arabic{enumi})}
  \tightlist
  \item
    Friendship from W1 - The results generalize pretty with the
    coefficients staying significant and positive across all classrooms
  \end{enumerate}
\item
  (2.i) Literacy receiver - The results are insignificant across all
  classrooms
\item
  (2.ii) Gender match - The results are significant across almost all
  classrooms and the coefficients are always positive. This suggest that
  our conclusion (that having the same gender is positively correlated
  with friendship in W2) should generalize.
\item
  (2.iii) HISEI sender - The results are almost always insignificant.
  Whenever they are significant, they are very close to zero. We can't
  say for sure whether the hypothesis is true or not.
\end{itemize}

\section{Task 3}\label{task-3}

\subsection{(1)}\label{section-5}

We can see that the network correlation coefficient is estimated to be
around \(0.85\) with a p-value of \(0\), which makes it significant.
Therefore, only taking W2 friendship into account, we can say that
students tend to mark friends people with similar literacy scores.

\subsection{(2) \& (3)}\label{section-6}

When we run the model with the additional variables, we get the
following estimates:

\begin{longtable}[]{@{}lrr@{}}
\caption{Extracted Results for Classroom 10}\tabularnewline
\toprule\noalign{}
Variable & Estimate & P\_Value \\
\midrule\noalign{}
\endfirsthead
\toprule\noalign{}
Variable & Estimate & P\_Value \\
\midrule\noalign{}
\endhead
\bottomrule\noalign{}
\endlastfoot
Gender (Female) & 20.8824892 & 0.0046561 \\
HISEI & 0.9727459 & 0.0000397 \\
Network Autocorrelation (rho) & -0.9395290 & 0.0282669 \\
\end{longtable}

We can see that all the coefficients are significant at the 5\%
significance level, therefore we can interpret them like so:

\begin{itemize}
\tightlist
\item
  Being a female (all other things staying the same) leads to a better
  literacy score. The estimated increase is around \(20.88\).
\item
  Having a higher HISEI score (all other things staying the same) leads
  to a better literacy score.
\item
  The network autocorrelation coefficient is now negative. In (1) it was
  estimated to be positive. This means that in the current model it is
  observed that people tend to make friends with others with different
  literacy score (higher or lower). The autocorrelation hypothesis is
  therefore suggested by the data.
\end{itemize}

\subsection{(4)}\label{section-7}

TODO

\subsection{(5)}\label{section-8}

\pandocbounded{\includegraphics[keepaspectratio]{solution_files/figure-latex/unnamed-chunk-4-1.pdf}}

We can see that our results do not look very generalizable and that
classroom 10 turns out to be an outlier. Observing the other classrooms,
we can see that almost every time the coefficients are non-significant,
which means that we cannot draw the same conclusions for other
classrooms.

\end{document}
